% 调用上级目录公共模板
\documentclass{../template/softeng}

\begin{document}

% 封面：文档名、项目名、版本、编写人、审核人、审批人、日期
\doccover{软件测试报告}
{待填写项目名称}
{V2026.07}

% 目录
\tableofcontents
\newpage

% 修订记录
\begin{revrecord}
V1.0 & 2026-07-12 & 测试编写人 & 测试负责人 & 初稿完成，搭建国标标准软件测试报告完整文档框架 \\
\hline
\end{revrecord}

\section{引言}
\subsection{编写目的}
\subsection{项目背景}
\subsection{术语与定义}
\subsection{参考资料}

\section{测试概要}
\subsection{测试项目概述}
\subsection{测试范围}
\subsection{测试环境}
\subsubsection{硬件环境}
\subsubsection{软件环境}
\subsubsection{网络环境}
\subsection{测试人员与周期}
\subsection{测试依据}

\section{测试内容与执行情况}
\subsection{功能测试}
\subsubsection{测试模块划分}
\subsubsection{测试用例执行统计}
\subsubsection{典型测试场景说明}
\subsection{性能测试}
\subsubsection{测试指标标准}
\subsubsection{压力/并发测试执行情况}
\subsection{兼容性测试}
\subsubsection{浏览器兼容测试}
\subsubsection{系统环境兼容测试}
\subsection{安全测试}
\subsubsection{权限校验测试}
\subsubsection{防注入、防XSS测试}
\subsection{界面易用性测试}

\section{测试用例执行结果统计}
\subsection{整体用例统计汇总}
\subsection{各模块用例通过率统计}
\subsection{缺陷分级统计（严重/一般/轻微/建议）}

\section{缺陷分析与问题汇总}
\subsection{缺陷分类统计}
\subsubsection{功能类缺陷}
\subsubsection{性能类缺陷}
\subsubsection{界面交互类缺陷}
\subsubsection{安全类缺陷}
\subsection{典型严重缺陷详情}
\subsection{未修复遗留缺陷说明}
\subsection{缺陷产生原因分析}

\section{测试结论}
\subsection{功能完整性结论}
\subsection{性能达标情况结论}
\subsection{系统安全与兼容性结论}
\subsection{上线/交付判定意见}

\section{建议与改进措施}
\subsection{程序代码优化建议}
\subsection{业务逻辑优化建议}
\subsection{后续迭代测试优化建议}

\section{附件}
\subsection{附件1：全部测试用例清单}
\subsection{附件2：缺陷完整记录表}
\subsection{附件3：性能测试监控截图与日志}
\subsection{附件4：测试执行原始日志}

% 插入测试截图/统计图表示例
\begin{figure}[htbp]
    \centering
    % \includegraphics[width=0.7\textwidth]{test_stat.png}
    \caption{测试缺陷分级统计饼图}
\end{figure}

% 示例代码/测试脚本块
\begin{lstlisting}[language=Python]
# 自动化测试脚本示例
def test_login(username, password):
    resp = request.post("/api/login", json={"username":username,"pwd":password})
    assert resp.status_code == 200
    return resp.json()
\end{lstlisting}

\end{document}